% SPDX-License-Identifier: MIT
% Copyright (c) 2026 Evgeny Baulin
% Optimization Methods, Seminar 09: one-page cheat sheet (Russian).
% Generated by checks/09. Linear Programming as the Foundation of Integer Optimization.ipynb from web/09/js/content-cheatsheet.js and content/students.json. Do not edit by hand.
% Compile with: latexmk -xelatex -interaction=nonstopmode -halt-on-error Cheatsheet_ru.tex
\RequirePackage{fix-cm}
\documentclass[10pt,a4paper]{article}

\usepackage{amsmath,amssymb}
\usepackage{fontspec}
\setmainfont{cmunrm.otf}[BoldFont=cmunbx.otf, ItalicFont=cmunti.otf, BoldItalicFont=cmunbi.otf]
\setsansfont{cmunss.otf}[BoldFont=cmunsx.otf, ItalicFont=cmunsi.otf, BoldItalicFont=cmunso.otf]
\setmonofont{cmuntt.otf}[BoldFont=cmuntb.otf, ItalicFont=cmunit.otf, BoldItalicFont=cmuntx.otf]
\usepackage{polyglossia}
\setdefaultlanguage{russian}
\setotherlanguage{english}
\newfontfamily\cyrillicfont{cmunrm.otf}[Ligatures=TeX, BoldFont=cmunbx.otf, ItalicFont=cmunti.otf, BoldItalicFont=cmunbi.otf]
\newfontfamily\cyrillicfontsf{cmunss.otf}[Ligatures=TeX, BoldFont=cmunsx.otf, ItalicFont=cmunsi.otf, BoldItalicFont=cmunso.otf]
\newfontfamily\cyrillicfonttt{cmuntt.otf}[BoldFont=cmuntb.otf, ItalicFont=cmunit.otf, BoldItalicFont=cmuntx.otf]
\PolyglossiaSetup{russian}{hyphenmins={2,3}}

\usepackage[a4paper,left=8mm,right=8mm,top=7mm,bottom=11mm,footskip=5.5mm]{geometry}
\usepackage{xcolor}
\usepackage{multicol}
\usepackage{enumitem}
\usepackage{ragged2e}
\usepackage{fancyhdr}
\usepackage[colorlinks=true,linkcolor=hseblue,urlcolor=hseblue]{hyperref}

\definecolor{hseblue}{RGB}{16,45,105}
\definecolor{hseaccent}{RGB}{198,40,40}
\definecolor{hsegray}{RGB}{108,114,126}
\definecolor{hselight}{RGB}{240,243,248}

\hypersetup{pdftitle={Шпаргалка к семинару 9: Линейное программирование и симплекс-метод}, pdfauthor={Евгений Баулин}}

\pagestyle{fancy}
\fancyhf{}
\renewcommand{\headrulewidth}{0pt}
\fancyfoot[L]{\scriptsize\color{hsegray}\textcopyright\ 2026 Евгений Баулин, \href{https://t.me/tarakan_tuc}{t.me/tarakan\_tuc}}
\fancyfoot[R]{\scriptsize\color{hsegray}Лицензия MIT}

\setlength{\parindent}{0pt}
\setlength{\parskip}{0pt}
\setlength{\columnsep}{4.5mm}
\setlength{\columnseprule}{0.3pt}
\renewcommand{\columnseprulecolor}{\color{hselight}}
\setlength{\emergencystretch}{2em}
\setlength{\RaggedRightRightskip}{0pt plus 4em}
\binoppenalty=2000
\relpenalty=500
\setlength{\abovedisplayskip}{2pt}
\setlength{\belowdisplayskip}{2pt}
\setlist[itemize]{leftmargin=7pt, labelsep=3pt, itemsep=1.2pt, topsep=1pt, parsep=0pt, partopsep=0pt,
  label={\color{hseaccent}\scriptsize\textbullet}}

\makeatletter
\newcommand{\cssec}[1]{\par\addvspace{3.5pt}\noindent
  \colorbox{hselight}{\parbox{\dimexpr\linewidth-2\fboxsep}{\RaggedRight\hyphenpenalty=10000\color{hseblue}\bfseries\strut #1}}%
  \par\nobreak\vspace{1pt}\@afterheading}
\makeatother
\clubpenalty=10000
\widowpenalty=10000

\begin{document}
\thispagestyle{fancy}

\noindent
\begin{minipage}[b]{0.6\textwidth}
  \RaggedRight\hyphenpenalty=10000
  {\color{hseaccent}\bfseries\small Шпаргалка к семинару 9}\\[1pt]
  {\color{hseblue}\bfseries\Large Линейное программирование\\ и симплекс-метод}
\end{minipage}\hfill
\begin{minipage}[b]{0.39\textwidth}
  \raggedleft\footnotesize
  {\bfseries Методы оптимизации}\\
  НИУ ВШЭ, факультет компьютерных наук, 2026--2027\\
  Лектор курса: Андрей Игнатов\\
  Автор: Евгений Баулин, Telegram \href{https://t.me/tarakan_tuc}{@tarakan\_tuc}
\end{minipage}

\vspace{2pt}
{\color{hseblue}\rule{\textwidth}{0.8pt}}
\vspace{-4pt}

\fontsize{9}{10.6}\selectfont
\begin{multicols}{3}
\RaggedRight


% ---- forms
\cssec{Формы записи}
\begin{itemize}
  \item $\max c^{\top}x$, $Ax\le b$, $x\ge0$; стандартная: $\min c^{\top}x$, $Ax=b$, $x\ge0$.
  \item Строка $\le$: $+s$; строка $\ge$: $-s$; свободная $x=x^+-x^-$; $b_i<0$: строка умножается на $-1$; $\max c^{\top}x=-\min(-c^{\top}x)$.
  \item Исходы: несовместна, неограниченна, есть оптимум.
\end{itemize}

% ---- basic
\cssec{Вершины}
\begin{itemize}
  \item Базис: $m$ независимых столбцов; $x_B=B^{-1}b$, остальные — ноль.
  \item Вершина $\Leftrightarrow$ базисное допустимое решение, $x_B\ge0$.
  \item Базисов не больше $\binom{n}{m}$. Оптимум, если есть, лежит в вершине.
\end{itemize}

% ---- step
\cssec{Шаг симплекс-метода (max)}
\begin{itemize}
  \item Последняя строка $z-c^{\top}x=0$; базис из дополнительных при $b\ge0$.
  \item Ввод: $r_j<0$ (Данциг: наиболее отрицательный; Бленд: наименьший индекс). Нет таких — оптимум.
  \item Вывод: наименьшее $b_i/a_{ie}$ при $a_{ie}>0$. Нет положительных — задача неограниченна.
  \item Шаг: столбец становится единичным, включая $z$-строку.
  \item Ноль под небазисной переменной: возможен другой оптимум.
\end{itemize}

% ---- phase
\cssec{Первая фаза}
\begin{itemize}
  \item Делаем $b\ge0$; строки без базисной переменной получают искусственные $a_i\ge0$.
  \item Минимизируем $w=\sum a_i$; сначала выражаем строку через небазисные.
  \item $w^*>0$: несовместна. $w^*=0$: выводим оставшиеся на нуле $a_i$, возвращаем $z$, выражаем через небазисные, продолжаем.
\end{itemize}

% ---- dual
\cssec{Двойственность}
\begin{itemize}
  \item $\max c^{\top}x$, $Ax\le b$, $x\ge0$ $\leftrightarrow$ $\min b^{\top}y$, $A^{\top}y\ge c$, $y\ge0$.
  \item Строка $=$ даёт свободное $y_i$; свободное $x_j$ — строку $=$.
  \item Слабая: $c^{\top}x\le b^{\top}y$. Сильная: в оптимуме равенство.
  \item Прямая неограниченна $\Rightarrow$ двойственная несовместна.
  \item Нежёсткость: $y_i(b-Ax)_i=0$, $x_j(A^{\top}y-c)_j=0$.
  \item $y^{\top}=c_B^{\top}B^{-1}$: $z$-строка под дополнительными.
\end{itemize}

% ---- sens
\cssec{Чувствительность}
\begin{itemize}
  \item $\partial p^*/\partial b_i=y_i^*$, пока держится базис; $p^*(b)$ вогнута и кусочно-линейна.
  \item Диапазон $b_i$: $x_B+\delta B^{-1}e_i\ge0$; $B^{-1}e_i$ — столбец $s_i$.
  \item Диапазон $c_j$: $z$-строка остаётся неотрицательной.
  \item Нежёсткая строка: цена ноль, но диапазон она ограничивать может.
\end{itemize}

% ---- breaks
\cssec{Сбои}
\begin{itemize}
  \item Вырожденность: базисная переменная — ноль, шаг может не сдвинуть точку.
  \item Зацикливание: базисы повторяются. Правило Бленда: наименьший индекс на входе и выходе.
  \item Худший случай — $2^n$ вершин (Кли — Минти); обычно шагов в несколько раз больше $m$.
  \item \texttt{linprog}: status 2 — несовместна, 3 — неограниченна; \texttt{ineqlin.marginals} — цены задачи min.
\end{itemize}

\end{multicols}
\end{document}
